\section{The Concept and the Failure of Warning}

\subsection{Aman, 'the Concept', and Fixed Assessments}
Israeli military intelligence, known as Aman, dominated national threat assessment in the years before 1973 and exuded confidence rooted in the triumph of 1967. Its analysts developed a governing assumption that came to be called ha-Konseptzia, or the Concept. According to this framework, Egypt would not launch war until it acquired long-range bombers capable of neutralizing Israeli airfields and achieved the air superiority needed to protect its forces. Until those conditions were met, an Egyptian offensive was judged improbable to the point of dismissal.

The Concept functioned less as a hypothesis than as an article of faith. Because Egypt had not obtained the requisite aircraft, Aman concluded that war was not imminent regardless of contrary signals. This rigidity reflected the broader hubris that followed 1967, when Israeli planners assumed permanent qualitative superiority. The assessment also rested on a misreading of Sadat, whom Israeli leaders underestimated as indecisive and unwilling to risk renewed defeat. The Concept thus encoded both technical reasoning and a flawed reading of an adversary's psychology.

This fixed assessment had organizational roots that magnified its danger. Aman held a near-monopoly on strategic evaluation, and dissenting voices struggled to penetrate the dominant narrative. A young analyst in the Sinai bureau who warned of imminent attack was overruled by superiors committed to the prevailing view. The institutional concentration of judgment meant that a single mistaken paradigm could shape the entire nation's posture, leaving little room for the competitive analysis that might have challenged comfortable assumptions before the crisis broke.

Senior figures including Moshe Dayan and the Aman director shared the conviction that no major Arab assault was likely in the autumn of 1973. Their confidence filtered down through the command structure, influencing decisions about mobilization, reserve readiness, and deployment along the canal. Because the political and military leadership trusted the Concept, they hesitated to call up reserves whose absence would prove costly. The assessment, in effect, governed not only analysis but the concrete operational choices that determined Israel's vulnerability on the morning of attack.

The persistence of the Concept illustrates how prior success can corrode strategic judgment. The very victories that established Israeli deterrence also bred an inflexibility that filtered evidence to fit expectation. Scholars continue to debate whether the failure stemmed primarily from this conceptual lock or from structural weaknesses in the warning system itself. Either way, the Concept stands as a classic case study in intelligence pathology, demonstrating how an organization can possess abundant information yet fail catastrophically to interpret it in time. \cite{siniver2013} \cite{dayan1976} \cite{bregman2002}

\subsection{Missed Indicators and the Agranat Commission}
In the days before 6 October a stream of warning indicators accumulated that, in hindsight, pointed unmistakably toward war. Egyptian and Syrian forces massed along the front lines under the cover of announced exercises. Soviet advisers and their families were suddenly evacuated, an alarming signal of impending conflict. Reconnaissance revealed unusual bridging equipment, forward artillery, and logistical preparations. Yet each indicator was reinterpreted through the lens of the Concept, which held that these movements reflected routine maneuvers rather than genuine offensive intent.

The failure to mobilize reserves in time compounded the analytical error. Israel's defense doctrine relied on rapid call-up of citizen-soldiers, but mobilization was delayed because leaders doubted the threat until the final hours. Only on the morning of 6 October did the picture clarify, and even then estimates of the attack's timing proved wrong, predicting late afternoon rather than the early-afternoon assault that actually came. The thin standing forces along the canal and the Golan were left to absorb the initial blow largely alone.

The trauma of the surprise prompted the Israeli government to establish a national commission of inquiry headed by Supreme Court President Shimon Agranat. The Agranat Commission examined the intelligence breakdown, the readiness failures, and the conduct of the war's opening days. Its investigation laid bare the dominance of the Concept and the dismissal of warning signs. The inquiry became a defining moment of national reckoning, forcing a society accustomed to military triumph to confront the limits of its own institutions and assumptions.

The commission's conclusions provoked enduring controversy. It assigned principal responsibility to the military command, recommending the dismissal of the Aman chief and the chief of staff, while largely sparing the political leadership from direct censure. This allocation struck many Israelis as unjust, since ministers had presided over the same complacency. Public anger nonetheless gathered force, and the perceived imbalance between military and political accountability shaped debate over civil-military relations and the proper limits of ministerial responsibility for national security.

Although the Agranat Commission formally exonerated Prime Minister Golda Meir and Defense Minister Moshe Dayan, the political consequences could not be contained. Mounting popular fury over the failures of October contributed to Meir's resignation in 1974 and intensified the criticism of Dayan that ended his ministerial career. The episode reshaped Israeli democracy by establishing a precedent of independent inquiry into security disasters. It demonstrated that even celebrated leaders could be held to account when institutions failed the public so visibly. \cite{siniver2013} \cite{rabinovich2004} \cite{dayan1976}

\begin{figure}[htbp]
\centering
\includegraphics[width=0.88\linewidth,height=0.58\textheight,keepaspectratio]{figures/ch02_web.jpg}
\caption{The Agranat Commission investigated the intelligence and command failures that left Israel unprepared on the morning of 6 October 1973.}
\label{fig:ch_the_agranat_commission_investiga}
\end{figure}

\bigskip
\begin{otherlanguage}{hebrew}
\subsection*{סיכום הפרק}

המודיעין הישראלי נצמד לקונספציה שמצרים לא תתקוף ללא עליונות אווירית ומפציצים ארוכי טווח. תפיסה קבועה זו גרמה לדחיית סימני האזהרה שהצטברו לפני המתקפה, וועדת אגרנט בחנה כישלון זה.

\end{otherlanguage}

\clearpage
